\documentclass[11pt]{article}

\usepackage[a4paper,margin=30mm]{geometry}
\usepackage[T1]{fontenc}
\usepackage{lmodern}
\usepackage{amsmath,amssymb,amsthm,mathtools}
\usepackage{array}
\usepackage{booktabs}
\usepackage{microtype}
\usepackage[hidelinks]{hyperref}

\allowdisplaybreaks
\setlength{\parindent}{0pt}
\setlength{\parskip}{0.65em}

\newtheorem{theorem}{Theorem}[section]
\newtheorem{lemma}[theorem]{Lemma}
\newtheorem{proposition}[theorem]{Proposition}
\newtheorem{corollary}[theorem]{Corollary}
\theoremstyle{remark}
\newtheorem{remark}[theorem]{Remark}

\newcommand{\Rbarp}{\overline{\mathbb R}_{+}}
\newcommand{\supp}{\operatorname{supp}}
\newcommand{\dd}{\,\mathrm{d}}
\newcommand{\ind}{\mathbf{1}}
\newcommand{\Aenergy}{\mathcal A}
\newcommand{\Senergy}{\mathcal S}
\newcommand{\Tenergy}{\mathcal T}
\newcommand{\Benergy}{\mathcal B}
\newcommand{\Venergy}{\mathcal V}
\newcommand{\Cube}{Q_d}
\newcommand{\Corners}{\mathcal C_d}

\title{Exact-Diameter Equivalence of Averaged and Singular Besov Energies
for Ahlfors-Regular Measures on the Cube}
\author{}
\date{}

\begin{document}

\maketitle

\begin{abstract}
Let a Borel measure be Ahlfors regular through the scale
\(R=\sqrt d\), and suppose that its support lies in the Euclidean unit
cube \([0,1]^d\).  We prove, with explicit positive finite structural
constants, the equivalence between an averaged, ball-normalized quadratic
energy and the corresponding singular Besov energy.  All energies take
values in the extended nonnegative reals, so the result requires no
finiteness assumption on the energies.  The exact cutoff \(R=\sqrt d\)
requires care because the balls are open: an open ball of radius equal to
the cube diameter need not contain the opposite corner.  We resolve this
endpoint by a quantitative corner-avoidance construction.  It produces a
positive-measure anchor set at a radius strictly below \(R\), which in turn
controls the long-range part of the singular energy.  The proof records the
exact constants and the order in which they are selected.  It is the
human-readable counterpart of a Lean 4 and Mathlib formalization.
\end{abstract}

\tableofcontents

\section{Setting and conventions}

Let
\[
  d\in\mathbb N,\qquad d>0,\qquad
  X=\mathbb R^d,\qquad
  \Cube=[0,1]^d,\qquad
  R=\sqrt d.
\]
The space \(X\) carries its Euclidean metric and Borel sigma-algebra.
Every ball \(B(x,r)\) in this article is an \emph{open} metric ball.

Let \(\nu\) be a Borel measure on \(X\), let
\[
  E=\supp\nu,
\]
and suppose
\[
  E\subseteq\Cube.
\]
Fix real parameters
\[
  \alpha>0,\qquad s>0,\qquad s<\frac{\alpha}{2},
  \qquad 0<c_A\leq C_A,
\]
and put
\[
  p=\alpha+2s.
\]
In particular, \(R>0\), \(C_A>0\), and \(p>0\).

The Ahlfors hypotheses are
\begin{equation*}
  c_A r^\alpha
  \leq \nu(B(x,r))
  \leq C_A r^\alpha
  \qquad
  (x\in E,\ 0<r\leq R).
  \tag{AR}
\end{equation*}

All nonnegative integrals are interpreted as extended nonnegative Lebesgue
integrals and take values in
\[
  \Rbarp=[0,\infty].
\]
We write
\[
  \iota:[0,\infty)\longrightarrow\Rbarp
\]
for the canonical embedding.  The extended inverse and product conventions
are
\[
  0^{-1}=\infty,\qquad
  \infty^{-1}=0,\qquad
  0\cdot\infty=\infty\cdot0=0.
\]
After the exact formal constants have been displayed, positive finite real
numbers will be identified with their images under \(\iota\).

Let \(u:X\to\mathbb R\) be Borel measurable and define
\[
  \Delta_u(x,y)=(u(x)-u(y))^2.
\]
For \(r\in\mathbb R\), set
\begin{align}
  D_r(x)
  &=
  \int_{B(x,r)}\Delta_u(x,y)\,\dd\nu(y),                                      \label{eq:ball-energy}\\
  L_r
  &=
  \int_XD_r(x)\,\dd\nu(x),                                                     \label{eq:local-energy}\\
  N_r
  &=
  \int_X\nu(B(x,r))^{-1}D_r(x)\,\dd\nu(x).                                    \label{eq:normalized-local-energy}
\end{align}
The averaged energy is
\begin{equation}
  \Aenergy_{s,R}(u)
  =
  \int_{(0,R]}r^{-1-2s}N_r\,\dd r,                                             \label{eq:averaged-energy}
\end{equation}
and the unnormalized scale energy is
\begin{equation}
  \Senergy_{\alpha,s,R}(u)
  =
  \int_{(0,R]}r^{-1-\alpha-2s}L_r\,\dd r
  =
  \int_{(0,R]}r^{-1-p}L_r\,\dd r.                                              \label{eq:scale-energy}
\end{equation}

The singular integrand is defined explicitly on the diagonal:
\[
  \Phi_{p,u}(x,y)
  =
  \begin{cases}
    0,&x=y,\\
    \Delta_u(x,y)\lvert x-y\rvert^{-p},&x\neq y.
  \end{cases}
\]
The singular Besov energy and global quadratic variation are
\begin{align}
  \Benergy_{\alpha,s}(u)
  &=
  \int_X\int_X\Phi_{p,u}(x,y)\,\dd\nu(y)\,\dd\nu(x),                            \label{eq:besov-energy}\\
  \Venergy(u)
  &=
  \int_X\int_X\Delta_u(x,y)\,\dd\nu(y)\,\dd\nu(x).                             \label{eq:global-variation}
\end{align}
The diagonal convention in \eqref{eq:besov-energy} is part of the
definition; the singular energy is not defined by silently evaluating a
product \(0\cdot\infty\).

\begin{remark}[The function space represented by the formal theorem]
The formal theorem quantifies over globally Borel measurable functions
\(u:X\to\mathbb R\).  It therefore applies, in particular, to any globally
measurable representative of an \(L^2(E,\nu)\) function.  It is not a
theorem stated on an \(L^2\) quotient type, and it does not formalize the
selection or extension of a representative defined only on \(E\).
\end{remark}

\section{Exact structural constants and the main theorem}

Let
\[
  \Corners=\{0,1\}^d
\]
be the set of cube corners.  Since \(\alpha>0\), one can choose, once and
for all and before choosing \(\nu\) or \(u\), a number \(\delta>0\) such
that
\begin{equation}
  2\delta\leq R                                                        \label{eq:delta-radius}
\end{equation}
and
\begin{equation}
  \sum_{v\in\Corners} C_A(2\delta)^\alpha
  <
  c_A R^\alpha.                                                        \label{eq:delta-small}
\end{equation}
Equivalently, the left side is \(2^d C_A(2\delta)^\alpha\).

Define the real geometric quantities
\begin{align}
  \rho_\delta
  &=
  \sqrt{d-\delta^2},                                                    \label{eq:rho}\\
  \varepsilon_\delta
  &=
  \frac{R-\rho_\delta}{2},                                              \label{eq:epsilon}\\
  \lambda_\delta
  &=
  \frac{R+\rho_\delta}{2R},                                             \label{eq:lambda}\\
  r_\delta
  &=
  \lambda_\delta R
  =
  \frac{R+\rho_\delta}{2}.                                              \label{eq:anchor-scale}
\end{align}
The exact constants used by the formal theorem are the following elements
of \(\Rbarp\):
\begin{align}
  m_-
  &=
  \iota\!\left(c_A\varepsilon_\delta^\alpha\right),                     \label{eq:mminus}\\
  m_+
  &=
  \iota\!\left(C_A R^\alpha\right),                                    \label{eq:mplus}\\
  K_\delta
  &=
  \iota\!\left(
    \frac{r_\delta^{-p}-R^{-p}}{p}
  \right)^{-1},                                                         \label{eq:tail-constant}\\
  G_\delta
  &=
  \left(\frac{4m_+}{m_-}\right)K_\delta,                                \label{eq:global-constant}\\
  D_\delta
  &=
  \left(
    \iota(p)+\iota(R)^{-p}G_\delta
  \right)\iota(C_A),                                                    \label{eq:reverse-constant}\\
  c
  &=
  D_\delta^{-1},                                                        \label{eq:lower-constant}\\
  C
  &=
  \iota(c_A)^{-1}\iota(p)^{-1}.                                        \label{eq:upper-constant}
\end{align}
Here \(\iota(R)^{-p}\) denotes the extended nonnegative real power.  Since
all of the relevant real quantities are positive and finite, it agrees
with \(\iota(R^{-p})\).

\begin{theorem}[Exact-diameter equivalence]                               \label{thm:main}
Fix \(d,\alpha,s,c_A,C_A\) as above, choose a structural \(\delta\)
satisfying \eqref{eq:delta-radius}--\eqref{eq:delta-small}, and define
\(c,C\) by \eqref{eq:rho}--\eqref{eq:upper-constant}.  Then
\[
  0<c<\infty,\qquad 0<C<\infty.
\]
For every Borel measure \(\nu\) on \(X\) and every globally Borel
measurable \(u:X\to\mathbb R\), if \(\nu\) satisfies \((\mathrm{AR})\) and
\(\supp\nu\subseteq\Cube\), then
\begin{equation}
  c\,\Benergy_{\alpha,s}(u)
  \leq
  \Aenergy_{s,R}(u)
  \leq
  C\,\Benergy_{\alpha,s}(u).                                             \label{eq:main-comparison}
\end{equation}
The constants are selected in the order
\[
  (d,\alpha,s,c_A,C_A)
  \longmapsto
  (\delta,c,C)
  \longmapsto
  (\nu,u).
\]
In particular, \(c\) and \(C\) do not depend on \(\nu\) or \(u\).
\end{theorem}

The remainder of the article proves Theorem~\ref{thm:main}.

\section{Measure-theoretic preliminaries}

\begin{lemma}[Automatic finiteness and sigma-finiteness]                  \label{lem:finite-measure}
Under the hypotheses of Theorem~\ref{thm:main}, the measure \(\nu\) is
finite.  Consequently it is sigma-finite.
\end{lemma}

\begin{proof}
The support \(E=\supp\nu\) is closed and lies in the compact cube
\(\Cube\), hence \(E\) is compact.  For every \(x\in E\), the upper
Ahlfors estimate at radius \(R\) gives
\[
  \nu(B(x,R))
  \leq C_A R^\alpha
  <\infty.
\]
The balls \(B(x,R)\), \(x\in E\), form an open cover of \(E\).  Compactness
gives a finite subcover, so finite subadditivity yields \(\nu(E)<\infty\).
The complement of the support is \(\nu\)-null in the second-countable
Euclidean space.  Therefore
\[
  \nu(X)=\nu(E)<\infty.
\]
A finite measure is sigma-finite.
\end{proof}

\begin{lemma}[Measurability of the energy kernels]                         \label{lem:measurability}
If \(u:X\to\mathbb R\) is Borel measurable, then all integrands in
\eqref{eq:ball-energy}--\eqref{eq:global-variation} are measurable.  The
maps \(x\mapsto D_r(x)\), \(r\mapsto L_r\), and \(r\mapsto N_r\) are
measurable wherever they are used.  All applications of Tonelli below are
therefore valid.
\end{lemma}

\begin{proof}
The map
\[
  (x,y)\longmapsto\Delta_u(x,y)
\]
is Borel measurable because subtraction and squaring are continuous.
The distance map is continuous, and hence
\[
  \{(r,x,y)\in\mathbb R\times X\times X:\lvert x-y\rvert<r\}
\]
is Borel.  Thus
\[
  (r,x,y)
  \longmapsto
  \ind_{\{\lvert x-y\rvert<r\}}\Delta_u(x,y)
\]
is a nonnegative measurable function.  Parameterized Tonelli, using the
sigma-finiteness from Lemma~\ref{lem:finite-measure}, shows that its
integrals in \(y\), and subsequently in \(x\), are measurable functions of
the remaining variables.  Applying the same argument to the indicator
alone shows that \((r,x)\mapsto\nu(B(x,r))\) is measurable.  Extended
inversion, multiplication, and real power are Borel operations on
\(\Rbarp\).  This proves the assertions for \(D_r\), \(L_r\), \(N_r\), and
the scale integrands.  The kernels in
\eqref{eq:besov-energy}--\eqref{eq:global-variation} are measurable by the
same closure properties, with the diagonal handled by a measurable
piecewise definition.
\end{proof}

\begin{lemma}[Atomlessness]                                                \label{lem:atomless}
Every singleton has zero \(\nu\)-measure.
\end{lemma}

\begin{proof}
First let \(x\in E\).  For every \(0<r\leq R\),
\[
  \nu(\{x\})
  \leq\nu(B(x,r))
  \leq C_A r^\alpha.
\]
Since \(\alpha>0\), the right side tends to zero as \(r\downarrow0\), and
therefore \(\nu(\{x\})=0\).  If \(x\notin E\), the definition of support
provides an open neighborhood \(U\) of \(x\) with \(\nu(U)=0\), whence
\(\nu(\{x\})=0\) as well.
\end{proof}

\begin{remark}
Lemma~\ref{lem:atomless} is independently useful, but the formal definition
of \(\Benergy_{\alpha,s}\) does not rely on it: the diagonal integrand is
set to zero explicitly.
\end{remark}

\section{Normalization by Ahlfors regularity}

\begin{lemma}[Fixed-scale normalization]                                  \label{lem:fixed-normalization}
For \(0<r\leq R\),
\[
  (C_A r^\alpha)^{-1}L_r
  \leq N_r
  \leq(c_A r^\alpha)^{-1}L_r.
\]
\end{lemma}

\begin{proof}
The support is conull.  Thus for almost every \(x\), \(x\in E\), and
\((\mathrm{AR})\) gives
\[
  c_A r^\alpha
  \leq\nu(B(x,r))
  \leq C_A r^\alpha.
\]
All three quantities are positive and finite.  Extended inversion reverses
the inequalities, so
\[
  (C_A r^\alpha)^{-1}
  \leq\nu(B(x,r))^{-1}
  \leq(c_A r^\alpha)^{-1}.
\]
Multiplication by \(D_r(x)\in\Rbarp\) preserves the order.  Integration in
\(x\) proves the result.
\end{proof}

\begin{lemma}[Integrated normalization]                                   \label{lem:integrated-normalization}
One has
\begin{equation}
  C_A^{-1}\Senergy_{\alpha,s,R}(u)
  \leq
  \Aenergy_{s,R}(u)
  \leq
  c_A^{-1}\Senergy_{\alpha,s,R}(u).                                      \label{eq:normalization-comparison}
\end{equation}
In particular,
\begin{equation}
  \Senergy_{\alpha,s,R}(u)
  \leq
  C_A\Aenergy_{s,R}(u).                                                   \label{eq:scale-by-averaged}
\end{equation}
\end{lemma}

\begin{proof}
For \(A>0\) and \(r>0\),
\[
  r^{-1-2s}(A r^\alpha)^{-1}
  =
  A^{-1}r^{-1-\alpha-2s}.
\]
Multiply the inequalities of Lemma~\ref{lem:fixed-normalization} by
\(r^{-1-2s}\), integrate over \((0,R]\), and use
\eqref{eq:averaged-energy}--\eqref{eq:scale-energy}.  This gives
\eqref{eq:normalization-comparison}.  Since \(C_A\) is positive and finite,
multiplication by \(C_A\) and cancellation of \(C_A C_A^{-1}=1\) yield
\eqref{eq:scale-by-averaged}.  The cancellation is valid in \(\Rbarp\)
because the cancelled factor is neither zero nor infinity.
\end{proof}

\section{The scale integral and its exact Tonelli kernel}

Define
\begin{equation}
  K_{p,R}(t)
  =
  \begin{cases}
    \dfrac{t^{-p}-R^{-p}}{p},&0<t\leq R,\\[2mm]
    0,&\text{otherwise}.
  \end{cases}                                                            \label{eq:truncated-kernel}
\end{equation}
The associated truncated double energy is
\begin{equation}
  \Tenergy_{p,R}(u)
  =
  \int_X\int_X
  K_{p,R}(\lvert x-y\rvert)\Delta_u(x,y)
  \,\dd\nu(y)\,\dd\nu(x).                                                \label{eq:truncated-energy}
\end{equation}

\begin{lemma}[Scalar scale integral]                                      \label{lem:scalar-scale}
If \(p>0\) and \(0<t\leq R\), then
\[
  \int_{(t,R]}r^{-1-p}\,\dd r
  =
  \frac{t^{-p}-R^{-p}}{p}.
\]
\end{lemma}

\begin{proof}
The function \(-p^{-1}r^{-p}\) is an antiderivative of
\(r^{-1-p}\) on \((0,\infty)\).  Evaluation at \(t\) and \(R\) gives the
formula.  The integrand is integrable on \((t,R]\) because the interval is
bounded away from zero.
\end{proof}

\begin{lemma}[Exact Tonelli identity]                                     \label{lem:tonelli}
For every measurable \(u:X\to\mathbb R\),
\begin{equation}
  \Senergy_{\alpha,s,R}(u)
  =
  \Tenergy_{p,R}(u).                                                       \label{eq:tonelli-identity}
\end{equation}
\end{lemma}

\begin{proof}
Expand \(L_r\) from \eqref{eq:local-energy} and
\eqref{eq:ball-energy}:
\[
  \Senergy_{\alpha,s,R}(u)
  =
  \int_{(0,R]}\int_X\int_X
  \ind_{\{\lvert x-y\rvert<r\}}
  r^{-1-p}\Delta_u(x,y)
  \,\dd\nu(y)\,\dd\nu(x)\,\dd r.
\]
The integrand is nonnegative and measurable by
Lemma~\ref{lem:measurability}, so Tonelli permits both changes of order.
Fix \(x,y\), and put \(t=\lvert x-y\rvert\).

If \(x\neq y\), then \(t>0\).  For \(0<t\leq R\), the allowed scales are
exactly \((t,R]\), and Lemma~\ref{lem:scalar-scale} gives
\[
  \int_{(0,R]}
  \ind_{\{t<r\}}r^{-1-p}\,\dd r
  =
  K_{p,R}(t).
\]
If \(t>R\), there is no allowed scale and both sides are zero.  If \(t=R\),
there is likewise no \(r\in(0,R]\) with \(t<r\); this is precisely why
\(K_{p,R}(R)=0\).

If \(x=y\), then \(\Delta_u(x,x)=0\), so the complete scale integrand is
identically zero.  We handle this case before factoring the scalar
integral: the unmultiplied integral of \(r^{-1-p}\) down to zero diverges,
and no informal infinity-times-zero calculation is used.

After integration in \(x\) and \(y\), the resulting expression is exactly
\eqref{eq:truncated-energy}.
\end{proof}

\section{Comparison with the full singular kernel}

\begin{lemma}[Pointwise kernel decomposition]                              \label{lem:kernel-decomposition}
For \(t>0\),
\begin{equation}
  t^{-p}
  \leq
  pK_{p,R}(t)+R^{-p}.                                                      \label{eq:kernel-upper}
\end{equation}
Equality holds whenever \(0<t\leq R\).
\end{lemma}

\begin{proof}
If \(0<t\leq R\), the definition gives
\[
  pK_{p,R}(t)+R^{-p}
  =
  t^{-p}-R^{-p}+R^{-p}
  =
  t^{-p}.
\]
If \(t>R\), then \(K_{p,R}(t)=0\).  Since \(p>0\), the function
\(t\mapsto t^{-p}\) is decreasing on \((0,\infty)\), and hence
\(t^{-p}\leq R^{-p}\).
\end{proof}

\begin{lemma}[Full-energy decomposition]                                  \label{lem:full-decomposition}
One has
\begin{equation}
  \Benergy_{\alpha,s}(u)
  \leq
  p\Tenergy_{p,R}(u)
  +
  R^{-p}\Venergy(u).                                                       \label{eq:full-decomposition}
\end{equation}
Moreover, equality holds if
\(\lvert x-y\rvert\leq R\) for \(\nu\otimes\nu\)-almost every pair
\((x,y)\).  In particular, equality holds under the cube-support
hypothesis.
\end{lemma}

\begin{proof}
On the diagonal, both \(\Phi_{p,u}(x,x)\) and
\(\Delta_u(x,x)\) are zero.  Off the diagonal, multiply
\eqref{eq:kernel-upper} by \(\Delta_u(x,y)\), and integrate twice.  Tonelli
and distributivity for nonnegative extended integrals give
\eqref{eq:full-decomposition}.

If almost every pair has distance at most \(R\), then the equality case of
Lemma~\ref{lem:kernel-decomposition} applies almost everywhere.  If
\(E\subseteq\Cube\), then
\[
  \lvert x-y\rvert^2
  =
  \sum_{j=1}^d(x_j-y_j)^2
  \leq d
\]
for \(x,y\in E\).  Since \(E\) is conull, this proves the final assertion.
Pairs at exact distance \(R\) are included: their truncated-kernel
contribution is zero, while their contribution appears in the
\(R^{-p}\Venergy(u)\) term.
\end{proof}

\begin{lemma}[Truncated energy is bounded by full energy]                  \label{lem:truncated-by-full}
One has
\begin{equation}
  \Tenergy_{p,R}(u)
  \leq
  p^{-1}\Benergy_{\alpha,s}(u).                                          \label{eq:truncated-by-full}
\end{equation}
\end{lemma}

\begin{proof}
For \(0<t\leq R\),
\[
  K_{p,R}(t)
  =
  \frac{t^{-p}-R^{-p}}p
  \leq
  p^{-1}t^{-p}.
\]
For \(t>R\), the left side is zero.  On the diagonal, the complete
integrands are both defined to be zero.  Multiplication by
\(\Delta_u(x,y)\) and two applications of monotonicity of the nonnegative
integral prove \eqref{eq:truncated-by-full}.
\end{proof}

\begin{corollary}[The easy comparison]                                    \label{cor:easy}
\[
  \Aenergy_{s,R}(u)
  \leq
  c_A^{-1}p^{-1}\Benergy_{\alpha,s}(u).
\]
\end{corollary}

\begin{proof}
By Lemmas~\ref{lem:integrated-normalization},
\ref{lem:tonelli}, and \ref{lem:truncated-by-full},
\[
  \Aenergy_{s,R}(u)
  \leq
  c_A^{-1}\Senergy_{\alpha,s,R}(u)
  =
  c_A^{-1}\Tenergy_{p,R}(u)
  \leq
  c_A^{-1}p^{-1}\Benergy_{\alpha,s}(u).
\]
\end{proof}

\section{A structural point away from every corner}

\begin{lemma}[Existence of the structural corner scale]                   \label{lem:corner-scale}
A number \(\delta>0\) satisfying
\eqref{eq:delta-radius}--\eqref{eq:delta-small} exists and depends only on
\(d,\alpha,c_A,C_A\).
\end{lemma}

\begin{proof}
Since \(\alpha>0\),
\[
  (2\delta)^\alpha\longrightarrow0
  \qquad(\delta\downarrow0).
\]
There are \(2^d\) corners, so
\[
  \sum_{v\in\Corners}C_A(2\delta)^\alpha
  =
  2^d C_A(2\delta)^\alpha
  \longrightarrow0.
\]
The target \(c_A R^\alpha\) is strictly positive.  Thus
\eqref{eq:delta-small} holds for all sufficiently small positive
\(\delta\).  Shrinking \(\delta\) further ensures \(2\delta\leq R\).
\end{proof}

\begin{lemma}[Corner-neighborhood estimate]                               \label{lem:corner-neighborhood}
For every corner \(v\in\Corners\),
\begin{equation}
  \nu(B(v,\delta))
  \leq
  C_A(2\delta)^\alpha.                                                     \label{eq:one-corner}
\end{equation}
Consequently, for
\[
  \mathcal N_\delta
  =
  \bigcup_{v\in\Corners}B(v,\delta),
\]
one has
\begin{equation}
  \nu(\mathcal N_\delta)
  \leq
  \sum_{v\in\Corners}C_A(2\delta)^\alpha
  <
  c_A R^\alpha.                                                           \label{eq:all-corners}
\end{equation}
\end{lemma}

\begin{proof}
Fix \(v\).  If \(\nu(B(v,\delta))=0\), then
\eqref{eq:one-corner} is immediate.  Otherwise the ball meets the support:
choose \(y_v\in B(v,\delta)\cap E\).  Indeed, the complement of \(E\) is
null, so a positive-measure Borel set cannot be contained in it.  If
\(z\in B(v,\delta)\), then
\[
  \lvert z-y_v\rvert
  \leq
  \lvert z-v\rvert+\lvert v-y_v\rvert
  <
  2\delta.
\]
Thus
\[
  B(v,\delta)\subseteq B(y_v,2\delta).
\]
Since \(y_v\in E\), \(2\delta>0\), and \(2\delta\leq R\), the upper
Ahlfors estimate gives
\[
  \nu(B(v,\delta))
  \leq
  \nu(B(y_v,2\delta))
  \leq
  C_A(2\delta)^\alpha.
\]
Finite subadditivity over the \(2^d\) corners, followed by
\eqref{eq:delta-small}, proves \eqref{eq:all-corners}.
\end{proof}

\begin{lemma}[Supported point uniformly away from the corners]            \label{lem:away-point}
Assume \(\nu\neq0\).  Then there exists \(q\in E\) such that
\begin{equation}
  \delta\leq\lvert q-v\rvert
  \qquad(v\in\Corners).                                                    \label{eq:away-corners}
\end{equation}
\end{lemma}

\begin{proof}
Since \(\nu\neq0\), the support is nonempty; choose \(x_0\in E\).  The
lower Ahlfors estimate at radius \(R\) gives
\begin{equation}
  c_A R^\alpha
  \leq
  \nu(B(x_0,R))
  \leq
  \nu(X).                                                                 \label{eq:mass-lower}
\end{equation}
If every supported point belonged to \(\mathcal N_\delta\), then conullity
of \(E\) would imply
\[
  \nu(X)=\nu(E)\leq\nu(\mathcal N_\delta).
\]
This contradicts \eqref{eq:all-corners} and
\eqref{eq:mass-lower}.  Hence some
\[
  q\in E\setminus\mathcal N_\delta
\]
exists.  Not belonging to the open ball \(B(v,\delta)\) is exactly
\(\lvert q-v\rvert\geq\delta\), proving
\eqref{eq:away-corners}.
\end{proof}

\section{Subdiametral geometry and the anchor set}

\begin{lemma}[Quantitative subdiametral radius]                            \label{lem:subdiametral}
Let \(q\in\Cube\) satisfy \eqref{eq:away-corners}.  Then
\[
  0\leq\rho_\delta<R
\]
and
\begin{equation}
  \lvert x-q\rvert
  \leq
  \rho_\delta
  \qquad(x\in\Cube).                                                       \label{eq:cube-in-subdiameter}
\end{equation}
\end{lemma}

\begin{proof}
Define a corner \(a\in\Corners\) coordinatewise by
\[
  a_j
  =
  \begin{cases}
    1,&q_j>\frac12,\\
    0,&q_j\leq\frac12.
  \end{cases}
\]
For every \(x\in\Cube\) and every coordinate \(j\),
\begin{equation}
  \lvert q_j-a_j\rvert^2
  +
  \lvert x_j-q_j\rvert^2
  \leq1.                                                                  \label{eq:coordinate-complement}
\end{equation}
Indeed, if \(q_j>1/2\), then \(a_j=1\) and
\[
  \lvert x_j-q_j\rvert\leq q_j,
  \qquad
  (1-q_j)^2+q_j^2\leq1.
\]
If \(q_j\leq1/2\), then \(a_j=0\) and
\[
  \lvert x_j-q_j\rvert\leq1-q_j,
  \qquad
  q_j^2+(1-q_j)^2\leq1.
\]
Summing \eqref{eq:coordinate-complement} gives
\[
  \lvert q-a\rvert^2+\lvert x-q\rvert^2\leq d.
\]
By \eqref{eq:away-corners}, \(\delta\leq\lvert q-a\rvert\), and therefore
\[
  \lvert x-q\rvert^2
  \leq
  d-\lvert q-a\rvert^2
  \leq
  d-\delta^2.
\]
This proves \eqref{eq:cube-in-subdiameter}.  It also shows
\(d-\delta^2\geq0\).  Since \(\delta>0\),
\[
  d-\delta^2<d,
\]
and monotonicity of the square root gives
\[
  \rho_\delta=\sqrt{d-\delta^2}<\sqrt d=R.
\]
\end{proof}

\begin{lemma}[Positive-measure anchor set below the diameter]             \label{lem:anchor-set}
Assume \(\nu\neq0\), and choose \(q\) as in
Lemma~\ref{lem:away-point}.  Let
\[
  S=B(q,\varepsilon_\delta)\cap E.
\]
Then \(S\) is measurable, \(S\subseteq E\), and
\begin{equation}
  m_-
  \leq
  \nu(S),                                                                 \label{eq:anchor-mass-lower}
\end{equation}
so in particular \(\nu(S)>0\).  Moreover,
\begin{equation}
  \lvert x-z\rvert<r_\delta
  \qquad(x\in E,\ z\in S),                                                \label{eq:anchor-distance}
\end{equation}
where
\[
  0<r_\delta<R.
\]
\end{lemma}

\begin{proof}
By Lemma~\ref{lem:subdiametral},
\[
  0\leq\rho_\delta<R.
\]
Hence
\[
  \varepsilon_\delta=\frac{R-\rho_\delta}{2}>0,
  \qquad
  r_\delta=\frac{R+\rho_\delta}{2}<R.
\]
Also \(\varepsilon_\delta\leq R\).  The open ball is Borel, the support is
closed, and therefore \(S\) is measurable.  Since \(E\) is conull,
\[
  \nu(S)=\nu(B(q,\varepsilon_\delta)).
\]
The point \(q\) belongs to \(E\), so the lower Ahlfors estimate gives
\[
  \nu(S)
  \geq
  c_A\varepsilon_\delta^\alpha
  =
  m_-.
\]
This quantity is strictly positive.

For \(x\in E\subseteq\Cube\) and \(z\in S\), use
Lemma~\ref{lem:subdiametral} and the strict ball membership of \(z\):
\[
  \lvert x-z\rvert
  \leq
  \lvert x-q\rvert+\lvert q-z\rvert
  <
  \rho_\delta+\varepsilon_\delta
  =
  \frac{R+\rho_\delta}{2}
  =
  r_\delta.
\]
This proves \eqref{eq:anchor-distance}.
\end{proof}

\begin{lemma}[Total-mass upper bound]                                      \label{lem:total-mass-upper}
Under the assumptions of Lemma~\ref{lem:anchor-set},
\begin{equation}
  \nu(X)\leq m_+.                                                          \label{eq:total-mass-upper}
\end{equation}
\end{lemma}

\begin{proof}
Because \(\nu(S)>0\), the set \(S\) is nonempty; choose \(z\in S\).  Then
\(z\in E\).  By \eqref{eq:anchor-distance} and \(r_\delta<R\),
\[
  E\subseteq B(z,R).
\]
This containment is strict at the metric level: every supported point has
distance less than \(r_\delta\), not merely at most \(R\).  Since \(E\) is
conull,
\[
  \nu(X)
  \leq
  \nu(B(z,R))
  \leq
  C_A R^\alpha
  =
  m_+,
\]
where the second inequality is the upper Ahlfors estimate at \(z\) and
radius \(R\).
\end{proof}

\begin{remark}[Why the corner construction is needed]
It is false in general that an open ball \(B(x,R)\), with
\(R=\sqrt d\), contains the whole cube: opposite corners are exactly
distance \(R\) apart.  Lemma~\ref{lem:anchor-set} replaces this invalid
endpoint argument by a structural scale \(r_\delta<R\).  The strict gap
simultaneously supplies an open-ball covering in
Lemma~\ref{lem:total-mass-upper} and a positive terminal interval of scales
below.
\end{remark}

\section{The anchor estimate for global variation}

\begin{lemma}[Three-point squared-difference inequality]                  \label{lem:three-point}
For all \(x,y,z\in X\),
\begin{equation}
  \Delta_u(x,y)
  \leq
  2\Delta_u(x,z)+2\Delta_u(z,y).                                         \label{eq:three-point}
\end{equation}
\end{lemma}

\begin{proof}
Write
\[
  u(x)-u(y)
  =
  (u(x)-u(z))+(u(z)-u(y)).
\]
The elementary inequality \((a+b)^2\leq2a^2+2b^2\) gives
\eqref{eq:three-point}.
\end{proof}

\begin{lemma}[Anchor control of global variation]                          \label{lem:anchor-variation}
Assume \(\nu\neq0\), and let \(S\) be the anchor set of
Lemma~\ref{lem:anchor-set}.  Then
\begin{equation}
  \Venergy(u)
  \leq
  \frac{4m_+}{m_-}\,L_{r_\delta}.                                        \label{eq:variation-by-local}
\end{equation}
\end{lemma}

\begin{proof}
Define
\[
  H(x)
  =
  \int_S\Delta_u(x,z)\,\dd\nu(z).
\]
Integrate \eqref{eq:three-point} over \(z\in S\).  Symmetry of
\(\Delta_u\) gives
\begin{equation}
  \nu(S)\Delta_u(x,y)
  \leq
  2H(x)+2H(y).                                                            \label{eq:point-anchor}
\end{equation}
Integrating in \(x\) and \(y\) and using Tonelli yields
\begin{align*}
  \nu(S)\Venergy(u)
  &\leq
  2\nu(X)\int_XH(x)\,\dd\nu(x)
  +
  2\nu(X)\int_XH(y)\,\dd\nu(y)\\
  &=
  4\nu(X)\int_XH(x)\,\dd\nu(x).
\end{align*}
For almost every \(x\), one has \(x\in E\).  By
\eqref{eq:anchor-distance}, \(S\subseteq B(x,r_\delta)\), and hence
\[
  H(x)
  \leq
  \int_{B(x,r_\delta)}\Delta_u(x,z)\,\dd\nu(z)
  =
  D_{r_\delta}(x).
\]
It follows that
\begin{equation}
  \nu(S)\Venergy(u)
  \leq
  4\nu(X)L_{r_\delta}.                                                     \label{eq:uncancelled-anchor}
\end{equation}
By \eqref{eq:anchor-mass-lower} and
\eqref{eq:total-mass-upper},
\[
  0<m_-\leq\nu(S)\leq\nu(X)\leq m_+<\infty.
\]
Thus division by \(\nu(S)\) is legitimate in \(\Rbarp\), and
\[
  \Venergy(u)
  \leq
  \frac{4\nu(X)}{\nu(S)}L_{r_\delta}
  \leq
  \frac{4m_+}{m_-}L_{r_\delta}.
\]
\end{proof}

\section{Recovering one local scale from the scale integral}

\begin{lemma}[Monotonicity of local energy]                               \label{lem:local-monotone}
The map \(r\mapsto L_r\) is nondecreasing.
\end{lemma}

\begin{proof}
If \(r\leq t\), then \(B(x,r)\subseteq B(x,t)\) for every \(x\).  Since
\(\Delta_u\geq0\),
\[
  D_r(x)\leq D_t(x).
\]
Integration in \(x\) gives \(L_r\leq L_t\).
\end{proof}

\begin{lemma}[Terminal scale interval]                                    \label{lem:terminal-scale}
Under the assumptions of Lemma~\ref{lem:anchor-set},
\begin{equation}
  L_{r_\delta}
  \leq
  K_\delta\Senergy_{\alpha,s,R}(u).                                      \label{eq:local-by-scale}
\end{equation}
\end{lemma}

\begin{proof}
Because \(0<r_\delta<R\) and \(p>0\),
\begin{equation}
  k_\delta
  :=
  \int_{(r_\delta,R]}r^{-1-p}\,\dd r
  =
  \frac{r_\delta^{-p}-R^{-p}}p
  >0.                                                                    \label{eq:tail-mass}
\end{equation}
For \(r\in(r_\delta,R]\), Lemma~\ref{lem:local-monotone} gives
\(L_{r_\delta}\leq L_r\).  Therefore
\[
  k_\delta L_{r_\delta}
  \leq
  \int_{(r_\delta,R]}r^{-1-p}L_r\,\dd r
  \leq
  \Senergy_{\alpha,s,R}(u).
\]
The number \(k_\delta\) is positive and finite, so extended cancellation
gives
\[
  L_{r_\delta}
  \leq
  k_\delta^{-1}\Senergy_{\alpha,s,R}(u)
  =
  K_\delta\Senergy_{\alpha,s,R}(u).
\]
\end{proof}

\begin{corollary}[Structural control of global variation]                 \label{cor:global-by-scale}
If \(\nu\neq0\), then
\begin{equation}
  \Venergy(u)
  \leq
  G_\delta\Senergy_{\alpha,s,R}(u)
  =
  G_\delta\Tenergy_{p,R}(u).                                              \label{eq:global-by-truncated}
\end{equation}
\end{corollary}

\begin{proof}
Combine Lemmas~\ref{lem:anchor-variation} and
\ref{lem:terminal-scale}:
\[
  \Venergy(u)
  \leq
  \frac{4m_+}{m_-}L_{r_\delta}
  \leq
  \frac{4m_+}{m_-}K_\delta
  \Senergy_{\alpha,s,R}(u).
\]
The coefficient is \(G_\delta\) by
\eqref{eq:global-constant}, and
Lemma~\ref{lem:tonelli} gives the final equality.
\end{proof}

\section{Proof of the main theorem}

\begin{lemma}[Positivity and finiteness of the constants]                 \label{lem:constant-positivity}
The constants in \eqref{eq:mminus}--\eqref{eq:upper-constant} satisfy
\[
  0<c<\infty,\qquad0<C<\infty.
\]
\end{lemma}

\begin{proof}
From \(d>0\), \(R=\sqrt d>0\).  From
\eqref{eq:delta-radius}, \(\delta\leq R/2<R\), and hence
\[
  0\leq\rho_\delta=\sqrt{d-\delta^2}<R.
\]
Therefore
\[
  \varepsilon_\delta>0,
  \qquad
  0<r_\delta<R.
\]
Since \(\alpha,c_A,C_A,R,\varepsilon_\delta\) are positive,
\[
  0<m_-<\infty,\qquad0<m_+<\infty.
\]
Because \(p>0\) and \(r_\delta<R\), the negative power is strictly
decreasing:
\[
  r_\delta^{-p}>R^{-p}.
\]
Thus the real number in parentheses in \eqref{eq:tail-constant} is
positive and finite, and consequently
\[
  0<K_\delta<\infty.
\]
It follows successively from
\eqref{eq:global-constant} and \eqref{eq:reverse-constant} that
\[
  0<G_\delta<\infty,
  \qquad
  0<D_\delta<\infty.
\]
The inverse of a positive finite extended nonnegative real is again positive
and finite, proving the claim for \(c=D_\delta^{-1}\).  The same argument
applied to \(c_A\) and \(p\) proves the claim for
\(C=c_A^{-1}p^{-1}\).
\end{proof}

\begin{proof}[Proof of Theorem~\ref{thm:main}]
The constants \(\delta,c,C\) have already been selected solely from the
structural data
\[
  d,\quad\alpha,\quad s,\quad c_A,\quad C_A.
\]
Lemma~\ref{lem:constant-positivity} proves their positivity and finiteness.
Now fix an arbitrary qualifying measure \(\nu\) and measurable function
\(u\).

The upper inequality follows for every \(\nu\), including the zero measure,
from Corollary~\ref{cor:easy}:
\[
  \Aenergy_{s,R}(u)
  \leq
  c_A^{-1}p^{-1}\Benergy_{\alpha,s}(u)
  =
  C\Benergy_{\alpha,s}(u).
\]

For the reverse inequality, first suppose \(\nu=0\).  Then every ball
energy, local energy, scale energy, truncated energy, global variation, and
singular energy is zero.  In the normalized local energy,
\(\nu(B(x,r))^{-1}=\infty\) and \(D_r(x)=0\); the stipulated
\(\Rbarp\) convention gives \(\infty\cdot0=0\).  Thus
\(\Aenergy_{s,R}(u)=0\), and the reverse inequality is \(0\leq0\).

Suppose now that \(\nu\neq0\).  Corollary~\ref{cor:global-by-scale} and the
full-energy decomposition give
\begin{align*}
  \Benergy_{\alpha,s}(u)
  &\leq
  p\Tenergy_{p,R}(u)
  +
  R^{-p}\Venergy(u)\\
  &\leq
  \left(p+R^{-p}G_\delta\right)
  \Tenergy_{p,R}(u).
\end{align*}
Using the Tonelli identity and the lower side of the normalization
comparison in its rearranged form \eqref{eq:scale-by-averaged},
\begin{align*}
  \Benergy_{\alpha,s}(u)
  &\leq
  \left(p+R^{-p}G_\delta\right)
  \Senergy_{\alpha,s,R}(u)\\
  &\leq
  \left(p+R^{-p}G_\delta\right)
  C_A\Aenergy_{s,R}(u)\\
  &=
  D_\delta\Aenergy_{s,R}(u).
\end{align*}
The coefficient \(D_\delta\) is positive and finite.  Hence extended
cancellation is valid, and
\[
  D_\delta^{-1}\Benergy_{\alpha,s}(u)
  \leq
  \Aenergy_{s,R}(u).
\]
Since \(c=D_\delta^{-1}\), this is the required reverse inequality.
\end{proof}

\begin{remark}[The unused exponent restriction]
The proof uses only \(p=\alpha+2s>0\), which follows from
\(\alpha>0\) and \(s>0\).  The restriction \(s<\alpha/2\) remains in
Theorem~\ref{thm:main} to match the formal target, but it is not needed for
the comparison proved here.
\end{remark}

\section{Endpoint and extended-valued safeguards}

For clarity, the proof uses the following safeguards explicitly.

\begin{enumerate}
  \item The scale domain is \((0,R]\), while the metric balls are open.
  Thus a pair at distance \(R\) is in no ball at an allowed scale.
  Accordingly \(K_{p,R}(R)=0\).

  \item The missing contribution at distance \(R\) is not discarded.  It
  appears exactly in the \(R^{-p}\Venergy(u)\) term of the full-energy
  decomposition.

  \item The diagonal is handled before the scalar scale integral is
  factored.  The scalar integral
  \[
    \int_0^R r^{-1-p}\,\dd r=\infty
  \]
  diverges, but the complete diagonal integrand is identically zero because
  \(\Delta_u(x,x)=0\).

  \item The cube diameter estimate gives only
  \(\lvert x-y\rvert\leq R\).  It does not imply that an open
  \(R\)-ball centered at an arbitrary cube point contains the cube.
  The corner-avoidance construction produces the strict anchor scale
  \(r_\delta<R\).

  \item The zero measure satisfies the formal Ahlfors predicate vacuously,
  because its support is empty.  It is therefore a genuine case of the
  theorem and is treated separately.

  \item No subtraction of extended energies is used.  The only subtraction
  occurs between finite positive real powers in
  \(K_{p,R}\) and \(k_\delta\), before embedding into \(\Rbarp\).

  \item Every cancellation or division in \(\Rbarp\) is made only after
  the relevant coefficient has been proved positive and finite.
\end{enumerate}

\appendix

\section{Correspondence with the Lean 4 formalization}

The compiled final declaration is
\[
  \texttt{BesovVerification.}
  \texttt{exists\_exactCube\_besovEnergy\_equivalence\_constants}.
\]
Its quantifier order is:
\[
  \forall d,\alpha,s,c_A,C_A,\quad
  \exists c,C,\quad
  \forall\nu,u.
\]
It assumes
\[
  d>0,\quad
  \alpha>0,\quad
  s>0,\quad
  s<\alpha/2,\quad
  c_A>0,\quad
  c_A\leq C_A,
\]
and concludes the two inequalities in
\eqref{eq:main-comparison}, with \(0<c,C<\infty\).

\begin{center}
\small
\begin{tabular}{@{}
  >{\raggedright\arraybackslash}p{0.39\textwidth}
  >{\raggedright\arraybackslash}p{0.53\textwidth}@{}}
\toprule
Lean source file & Mathematical role in this article\\
\midrule
\texttt{Definitions.lean}
&
Ambient Euclidean space, unit cube, Ahlfors predicate, ball, local,
normalized, averaged, singular, and global energies.\\

\texttt{Measurability.lean}
&
Parameterized measurability for moving balls, ball measures, local
energies, and normalized scale integrands; supports
Lemma~\ref{lem:measurability}.\\

\texttt{Ahlfors.lean}
&
Atomlessness from the upper Ahlfors estimate; corresponds to
Lemma~\ref{lem:atomless}.\\

\texttt{CoreLemmas.lean}
&
The three-point squared-difference inequality and monotonicity of local
energy; Lemmas~\ref{lem:three-point} and \ref{lem:local-monotone}.\\

\texttt{ScaleIntegral.lean}
&
The real and extended-nonnegative scalar integral in
Lemma~\ref{lem:scalar-scale}.\\

\texttt{Normalization.lean}
&
Fixed-scale and integrated Ahlfors normalization;
Lemmas~\ref{lem:fixed-normalization} and
\ref{lem:integrated-normalization}.\\

\texttt{EnergyIdentity.lean}
&
Measurable triple kernel, diagonal handling, Tonelli swaps, and the exact
identity \eqref{eq:tonelli-identity}.\\

\texttt{MainTheorem.lean}
&
Comparison of the truncated kernel with the singular kernel and the easy
upper inequality; Lemma~\ref{lem:truncated-by-full} and
Corollary~\ref{cor:easy}.\\

\texttt{FullEnergyDecomposition.lean}
&
The pointwise singular-kernel identity and the global-variation tail;
Lemmas~\ref{lem:kernel-decomposition} and
\ref{lem:full-decomposition}.\\

\texttt{CubeAnchor.lean}
&
Corner neighborhoods, structural corner-scale existence, compact-support
finiteness, corner avoidance, subdiametral geometry, and quantitative
anchor balls; the compact-support, corner-avoidance, and anchor-geometry
arguments above.\\

\texttt{FixedCornerScaleAnchor.lean}
&
Packages the anchor construction with \(\delta\) fixed before the measure,
which is essential for the final quantifier order.\\

\texttt{AnchorEstimate.lean}
&
Integration of the three-point inequality over the anchor set;
Lemma~\ref{lem:anchor-variation}.\\

\texttt{StructuralAnchorEstimate.lean}
&
Replaces the actual anchor and total masses by the uniform bounds
\(m_-\) and \(m_+\).\\

\texttt{TailEstimate.lean}
&
Recovers \(L_{r_\delta}\) from the terminal scale interval;
Lemma~\ref{lem:terminal-scale}.\\

\texttt{AnchorTailComparison.lean}
&
Composes anchor, terminal-scale, and normalization estimates.\\

\texttt{ExactCubeConstants.lean}
&
Defines \(\rho_\delta,\varepsilon_\delta,\lambda_\delta,r_\delta\) and the
exact constants \eqref{eq:mminus}--\eqref{eq:upper-constant}, and proves
their positivity and finiteness.\\

\texttt{ExactCubeMainTheorem.lean}
&
Handles the zero measure, assembles both inequalities, and establishes the
final structural quantifier order.\\
\bottomrule
\end{tabular}
\end{center}

\section{Pinned verification environment}

The formal development was compiled in the following pinned environment:
\begin{center}
\begin{tabular}{@{}ll@{}}
\toprule
Component & Version\\
\midrule
Lean & \(4.32.1\)\\
Lean commit & \texttt{f054605aea4b840552cca2e725580bffd1e1b704}\\
Lake & \(5.0.0\)\\
Elan & \(4.2.3\)\\
Toolchain & \texttt{leanprover/lean4:v4.32.1}\\
Mathlib commit & \texttt{520045ab14e26149ee970e2e617ca04b09bde5d6}\\
\bottomrule
\end{tabular}
\end{center}

The final theorem's reported logical dependencies are
\[
  \texttt{propext},\qquad
  \texttt{Classical.choice},\qquad
  \texttt{Quot.sound}.
\]
These are standard Lean axioms used by Mathlib.  The checked project source
contains no unfinished proof term and introduces no custom axiom asserting
the mathematical result.

\end{document}
